% Compiles paper/hyperelliptic-candidates.tex on its own, with the macros and theorem environments of the
% manuscript (paper/manuscript/arXiv_v1.tex) it is written for.  pdflatex twice.
\documentclass[11pt]{amsart}
\usepackage[margin=1.1in]{geometry}
\usepackage{amsmath,amssymb,amsthm,mathtools,tabularx,array,xcolor}
\usepackage[colorlinks=true,linkcolor=blue,citecolor=blue,urlcolor=blue]{hyperref}
\newcommand{\todo}[1]{\textcolor{red}{[#1]}}
\newcommand{\Q}{\mathbb{Q}}
\newcommand{\Z}{\mathbb{Z}}
\newcommand{\F}{\mathbb{F}}
\DeclareMathOperator{\Jac}{Jac}
\DeclareMathOperator{\Pic}{Pic}
\newcommand{\Curve}[2]{X_0({#1},{#2})}
\newcommand{\StarCurve}[2]{X_0^*({#1},{#2})}
\newcommand{\ALfull}[2]{W_{{#1}{#2}}}
\theoremstyle{plain}
\newtheorem{theorem}{Theorem}[section]
\newtheorem{proposition}[theorem]{Proposition}
\newtheorem{lemma}[theorem]{Lemma}
\newtheorem{corollary}[theorem]{Corollary}
\theoremstyle{definition}
\newtheorem{remark}[theorem]{Remark}
% labels of the manuscript sections referred to, so that the references resolve here
\begin{document}
\title{The remaining candidates}
\author{}
\date{\today}
\maketitle
\setcounter{section}{8}
\makeatletter
\newcommand{\fakelabel}[2]{\@namedef{r@#1}{{#2}{}{}{}{}}}
\makeatother
\fakelabel{sec:results}{8}\fakelabel{sec:fixedpointsAL}{3.2}\fakelabel{sec:trace}{4}\fakelabel{sec:specialfiber}{8.4}\fakelabel{sec:limits}{9.2}
% The remaining candidates: a section for the manuscript (paper/manuscript/arXiv_v1.tex), written in its
% conventions (\Curve{D}{N}, \StarCurve{D}{N}, \ALfull{D}{N}; theorem environments numbered by equation) and meant
% to be \input after Section \ref{sec:results}.  Every general statement is proved in one of the notes under
% paper/ (cited by file and printed numbering); every per-curve claim is a test in tests/.  Numbers marked
% \todo{} must be re-read from the current manuscript before the section is merged into it.
% Compile alone with paper/hyperelliptic-candidates-standalone.tex.

\section{The remaining candidates}\label{sec:candidates}

The filters of Section~\ref{sec:results} leave a list of quotients $\Curve{D}{N}/W$ of genus at least three
whose hyperellipticity they neither establish nor exclude. This section describes what decides the question for
the curves on that list that are hyperelliptic. The answer has a uniform shape. A numerical computation of the
differentials of $\Curve{D}{N}/W$ says which of the open quotients look hyperelliptic and, for those, which pairs
of CM points look like fibres of the hyperelliptic map. Each such guess is a statement about divisor classes on a
quotient curve: that a certain degree-zero combination of CM divisors is principal. The torsion part of that
statement is settled by Borcherds products or, equivalently, by a linear-algebra condition on cusp forms of
weight $\tfrac32$; the torsion itself is settled by reduction at the primes of~$D$, where the quotient is a
Mumford curve whose special fibre, and the reduction of its CM points to it, are known explicitly. No equation of
the curve is needed at any point, and for most of the curves none is produced.

\subsection{Where the list stands}\label{sec:candidates-list}

\todo{Reconcile with the current manuscript: the filters added after the first version (the dual graph of the
special fibre at a prime of $D$, its isometric involutions, and the isogeny classes it determines) bring the
open list from $784$ to $61$; the count below refers to those $61$.}

Among the open quotients, a numerical criterion singles out $34$ as hyperelliptic. Let $Y=\Curve{D}{N}/W$ have
genus $g\ge3$ with a basis $u_1,\dots,u_g$ of holomorphic differentials, evaluated at sample points of the upper
half plane through Nelson's explicit form of the Shimizu lift (\texttt{paper/nelson-method.tex}, Section~1): the
weight-two form on $\Curve{D}{N}$ attached to a newform $f$ of level $DN$ is a sum over a positive-definite
quaternary lattice of Bessel terms weighted by the coefficients of $f$, so a basis of differentials of any
quotient is computable without a fundamental domain, the $W$-invariant forms being the $W$-fixed newforms of
level $DN$ together with the old classes whose Atkin--Lehner signs multiply to one on $W$
(\texttt{paper/nelson-method.tex}, Proposition~2.1). Noether's theorem then separates the two cases: the
canonical image of $Y$ lies on $\binom{g-1}{2}$ independent quadrics if $Y$ is hyperelliptic and on
$\binom{g-2}{2}$ otherwise, one quadric against none in genus three, three against one in genus four, and the
number is read off the singular values of the matrix of quadratic monomials at the sample points. The verdict is
numerical; it is accepted only when it is stable under the cutoff of the lattice sum and the working precision,
and it is used only to direct the exact arguments below. The same computation names, for each hyperelliptic
candidate, the pairs of CM points whose canonical images coincide, which are the fibres of the hyperelliptic map
through CM points, and these are the divisors the exact arguments prescribe. (At a CM point fixed by some
$w_m\in W$ every pulled-back differential vanishes, so the coordinate is $0/0$ there and its value is taken from
the derivatives; \texttt{paper/nelson-method.tex}, Section~7.)

Of the $34$, one is not hyperelliptic: $\Curve{141}{2}/\{1,6,94,141\}$ is a double cover of an elliptic curve
whose possible covers are enumerated exactly in Section~\ref{sec:candidates-heights} and none is hyperelliptic.
The other $33$ fall into three shapes according to the involutions that survive on the quotient.

\begin{enumerate}
\item \emph{Genus three over genus one.} $W$ has index two in a group $W'$ for which $E=\Curve{D}{N}/W'$ has
genus one; $Y\to E$ is a double cover branched at four points. Twenty cases.
\item \emph{Genus four over genus two.} $W$ has index two in $\ALfull{D}{N}$ and the star
$C=\StarCurve{D}{N}$ has genus two; $Y\to C$ is a double cover branched at two points. Six cases.
\item \emph{Genus three with the full group.} $Y=\StarCurve{D}{N}$ itself, of genus three, with no residual
involution. Seven cases.
\end{enumerate}

\subsection{Three criteria}\label{sec:candidates-criteria}

In each shape the branch locus of the double cover is known exactly. For $W\subset W'$ of index two, a point of
$\Curve{D}{N}/W'$ is a branch point of $\Curve{D}{N}/W\to\Curve{D}{N}/W'$ exactly when its $W'$-orbit on
$\Curve{D}{N}$ is a single $W$-orbit, that is, when the orbit consists of fixed points of some $w_m$ with
$m\in W'\setminus W$; by Ogg these are the CM points of discriminant $-4m$ and, for $m\equiv3\pmod4$, $-m$ (and
$-4,-8$ for $m=2$), counted by the formula of Section~\ref{sec:fixedpointsAL}, with the degree of the branch
divisor fixed by Riemann--Hurwitz (\texttt{paper/heegner-height-models.tex}, Section~2). The fields of definition
of the branch points follow from Shimura reciprocity.

\begin{proposition}[genus three over genus one]\label{prop:bielliptic}
Let $Y$ be a curve of genus three with an involution whose quotient $E$ has genus one, branched at
$B_1,\dots,B_4$. Then $Y$ is hyperelliptic if and only if the branch points split into two pairs with
$B_1+B_2\sim B_3+B_4$ on $E$.
\end{proposition}

The hyperelliptic involution is central, so it commutes with the covering involution and descends to $E$, whose
quotient is then a line; the degree-two map so obtained on $E$ has the fixed points of the covering involution in
two of its fibres. Conversely a function $x$ on $E$ with divisor $B_1+B_2-B_3-B_4$ pulls back to a square on $Y$
(\texttt{paper/hyperelliptic-15-38.tex}, Section~1). Over $\Q$ the double covers of $E$ branched exactly at a
given rational divisor form a torsor under $E(\Q)[2]$, up to constant twists, and only one of them is the
hyperelliptic one; which cover $Y$ is must be decided separately (Section~\ref{sec:candidates-torsion}).

\begin{proposition}[genus four over genus two]\label{prop:normalform}
Let $\pi\colon Y\to C$ be a double cover over $\Q$ of a curve of genus two by a curve of genus four, branched at
$P+Q$ with $P,Q\in C(\Q)$. Then $Y$ is hyperelliptic if and only if $P+Q$ is a fibre of the hyperelliptic map
of $C$; in that case $C$ has a rational Weierstrass point, $C\colon y^2=h(u)$ with $h$ a quintic, $h(0)\neq0$,
the fibre at $u=0$, and $Y$ is $y^2=h(t^2/c)$ for some $c\in\Q^\times$.
\end{proposition}

This is Lemma~8 of \texttt{paper/star-2002-proof.tex}. Since two distinct effective divisors of degree two on a
curve of genus two that are linearly equivalent are both canonical, the condition is that the branch pair be
linearly equivalent to any one other fibre.

\begin{proposition}[genus three with the full group]\label{prop:g12}
Let $Y$ be a curve of genus three. If $D_1\neq D_2$ are effective divisors of degree two with $D_1\sim D_2$, then
$|D_1|$ is a $g^1_2$ and $Y$ is hyperelliptic.
\end{proposition}

In all three shapes, then, hyperellipticity is the statement that a specific degree-zero divisor supported on CM
points, a difference of two candidate fibres, is principal. The rest of the section is about proving such a
statement.

\subsection{Relations between CM divisors}\label{sec:candidates-relations}

Let $X^*=\Curve{D}{N}/W$ be the quotient on which the relation is to be proved, $M=4\cdot\mathrm{oddpart}(DN)$,
and for $n\ge1$ let
\[
Z(n)=\sum_{r\ge1,\ d=-4n/r^2}\frac1{e_d}\,P_d
\]
be the Heegner divisor of index $n$ on $X^*$: $P_d$ is the sum of the CM points of discriminant $d$ on $X^*$ and
$e_d$ the order of their stabiliser (Guo--Yang, Lemma~25). A weakly holomorphic form of weight $\tfrac12$ on
$\Gamma_0(M)$ with the theta multiplier and principal part $\sum_n c(-n)q^{-n}$ at $\infty$ has a Borcherds lift
to $X^*$ with divisor $\sum_n c(-n)Z(n)$, and of weight zero exactly when this divisor has degree zero. Two
consequences are used.

\begin{lemma}\label{lem:borcherds-torsion}
If the principal part is integral, the class of the divisor of the lift is torsion in $\Jac(X^*)(\Q)$; with a
rational principal part the same holds for the divisor it prescribes.
\end{lemma}

This is Borcherds's theorem that the lift has a character of finite order
(\texttt{paper/torsion-counting-hyperelliptic.tex}, Lemma~1 and Corollary~2). When the degree-zero CM divisors
generate a torsion-free subgroup of the Jacobian, the character is trivial and the lift is a function
(\texttt{paper/character-and-reduction-types.tex}, Proposition~3.1 and Corollary~3.2); in general its order
divides the exponent of the rational torsion, which is why the torsion must be handled on its own.

\begin{lemma}\label{lem:dual-certificate}
A weakly holomorphic form of weight $\tfrac12$ on $\Gamma_0(M)$ with the theta multiplier, holomorphic away
from $\infty$ and with a prescribed principal part $\sum_n c(-n)q^{-n}$, exists if and only if
$\sum_n c(-n)a_g(n)=0$ for every cusp form $g$ of weight $\tfrac32$ on $\Gamma_0(M)$ with the conjugate
multiplier.
\end{lemma}

The product of the two forms is a differential on $X_0(M)$ with its only pole at $\infty$, so the condition is
the vanishing of a residue, and Serre duality makes it sufficient (\texttt{paper/hyperelliptic-6-379.tex},
Proposition~3.1). The lemma turns the existence of a Borcherds form with a given divisor into linear algebra:
the space $S_{3/2}(\Gamma_0(M))$ is computed once, to a precision at least the largest index the divisor
involves, and each prescribed divisor is a linear system in the $c(-n)$. Where a search for the form itself by
eta quotients (Section~\ref{sec:limits}) does not finish, the certificate decides in seconds.

The same duality gives a test for the torsion statement that needs no form at all.

\begin{proposition}[the dual test]\label{prop:dual-test}
Let $\lambda$ be the Hodge class in $\Pic(X^*)_\Q$, and for a linear form $\varphi$ on $\Pic(X^*)_\Q$ with
$\varphi(\lambda)=0$ put $\Phi_\varphi=\sum_{n\ge1}\varphi(Z(n))q^n$. Then $\Phi_\varphi$ is a cusp form of
weight $\tfrac32$ on $\Gamma_0(M)$ with the multiplier of $\theta^3$, vanishing at every index $n$ with
$Z(n)=0$. Consequently, if $\Delta=\sum_n c(-n)Z(n)$ is a degree-zero CM divisor and
$\sum_n c(-n)a_g(n)=0$ for every $g$ in the space $\mathcal W$ of such forms, then $\Delta$ is torsion in
$\Jac(X^*)(\Q)$.
\end{proposition}

(\texttt{paper/hyperelliptic-15-38.tex}, Proposition~2.1; \texttt{paper/hyperelliptic-21-58.tex},
Proposition~2.1.) The space $\mathcal W$ is cut out of $S_{3/2}(\Gamma_0(M))$ by the vanishing conditions, and
it is sharp when $2\mid D$: its dimension is the rank of $\Jac(X^*)(\Q)$ in every case computed, and the Hecke
operators $T(p^2)$ act on it through the traces of Frobenius of $X^*$. Since the forms $\Phi_\varphi$ lie in
$\mathcal W$ whatever its dimension, orthogonality to all of $\mathcal W$ proves torsion unconditionally; the
converse, that a divisor not orthogonal to $\mathcal W$ is not torsion, needs $\mathcal W$ to be spanned by the
$\Phi_\varphi$, and is used only for controls. The precision of the computation must reach the index $|d|/4$ of
every discriminant in the divisor; the cost is in the level, not the precision.

When $2\parallel N$ the dictionary is coarser: an odd discriminant $d$ and $4d$ are reached from the same
indices and both carry CM points, so $P_d$ and $P_{4d}$ cannot be prescribed separately, and almost every index
carries CM points, so $\mathcal W$ is twice too large. Both are repaired by passing to the curve of level $N/2$:
with $T$ the correspondence from $X^*$ to the star of level $N/2$, $T\circ Z=(U(4)+2)H$ for the Heegner series
$H$ of the lower level, $T$ is injective on $\Jac(X^*)_\Q$ when the Jacobian is old at $2$, and the test is run
against $(U(4)+2)$ applied to the Kohnen plus-space forms vanishing where $H$ vanishes
(\texttt{paper/hyperelliptic-21-58.tex}, Lemmas~3.1--3.3, Proposition~3.4 and Theorem~3.5); an odd-discriminant
fibre is prescribed directly on the lower curve (\texttt{paper/hyperelliptic-39-34.tex}, Proposition~2.1).

Finally, torsion relations are turned into principal ones by counting.

\begin{theorem}[counting]\label{thm:counting}
Let $t$ be a multiple of $\#\Jac(X^*)(\Q)_{\mathrm{tors}}$, and let $D_1,\dots,D_r$ be pairwise distinct
effective divisors of degree two on $X^*$, defined over $\Q$, such that $D_i-D_j$ is torsion for all $i\neq j$.
If $r>t$, then two of the $D_i$ are linearly equivalent and $X^*$ is hyperelliptic.
\end{theorem}

(\texttt{paper/torsion-counting-hyperelliptic.tex}, Theorem~3.) The bound $t$ comes from the trace formula: the
rational torsion injects into $\Jac(X^*)(\F_p)$ for every odd prime $p$ of good reduction, so its order divides
the greatest common divisor of the orders $\#\Jac(X^*)(\F_p)$, which are the values at $1$ of the zeta numerators
computed from point counts (Section~\ref{sec:trace}); no model is needed. When $t=1$ a single torsion relation
between two fibres suffices, and the test of Proposition~\ref{prop:dual-test} proves hyperellipticity by itself.
The divisors $D_i$ beyond those named by the numerical screen are found by the same test run as a search over
all CM divisors of degree at most two.

\subsection{The torsion part}\label{sec:candidates-torsion}

When the torsion bound is not $1$ and not enough fibres are available, or when the relation is a pairing on an
elliptic curve with rational torsion, the torsion part of the relation has to be decided, and this is done by
reduction at a prime $p\mid D$. There $\Curve{D}{N}$ has the \v{C}erednik--Drinfeld model, and $Y=\Curve{D}{N}/W$
inherits a semistable model whose special fibre is a union of rational curves meeting at ordinary double points.
Its dual graph is the quotient by $W$ of the graph of Section~\ref{sec:specialfiber}: vertices the two copies of
the left ideal classes of an Eichler order of level $N$ in the definite algebra of discriminant $D/p$, edges the
classes of the level-$Np$ suborder, with the thickness of a double point the order of the unit group of its
order modulo $\pm1$; under the quotient a double point whose two branches are swapped by an involution of $W$
divisible by $p$ becomes a smooth point, and a surviving double point has its thickness multiplied by the order
of its stabiliser.

\begin{proposition}\label{prop:component-group}
The component group $\Phi_p$ of the N\'eron model of $\Jac(Y)$ at $p$ is the critical group of the graph obtained
from the quotient graph by replacing each edge of thickness $t$ by a chain of $t$ edges: the torsion of the
cokernel of its Laplacian.
\end{proposition}

\begin{proposition}\label{prop:cm-reduction}
A CM point of discriminant $d$ on $Y$, $p$ unsplit in $\Q(\sqrt d)$, reduces to a double point if $p$ ramifies in
$\Q(\sqrt d)$ and to a smooth point of a component if $p$ is inert; the double point, or the component, is one
whose Eichler order admits an optimal embedding of the order of discriminant $d$, and the number of CM points of
$\Curve{D}{N}$ reducing to it is the number of such embeddings up to conjugation by the units of that order.
\end{proposition}

The first is Raynaud's theorem applied to the minimal regular model (\texttt{paper/component-groups.tex},
Proposition~1.1); the second is Molina's description of the specialisation of Heegner points, Theorems~5.3,
5.4 and 5.8 of \emph{Ribet bimodules and the specialization of Heegner points}, which rests on Ribet's
\emph{Bimodules and abelian surfaces} and carries no restriction on $p$. Together they make the class in
$\Phi_p$ of any degree-zero combination of CM divisors of $Y$ computable by counting optimal embeddings into
definite orders. A nonzero class shows the divisor is not principal. For the converse one needs to know that
the torsion in question is seen by $\Phi_p$.

\begin{proposition}\label{prop:torsion-injects}
Suppose the Hall divisor $p$ itself lies in $W$. Then the identity component of the special fibre of the N\'eron
model of $\Jac(Y)$ at $p$ is a split torus, and for every prime $\ell$ not dividing $p-1$, with $\ell=p$
allowed when $p\ge3$, the $\ell$-primary part of $\Jac(Y)(\Q)_{\mathrm{tors}}$ injects into $\Phi_p$. In
particular all odd torsion injects into $\Phi_2$, and the $p$-part into $\Phi_p$ for every odd $p\mid D$.
\end{proposition}

Frobenius acts on the \v{C}erednik--Drinfeld graph as $w_p$, hence trivially on the quotient graph when
$w_p\in W$, which splits the torus; its $\F_p$-points have order $(p-1)^g$, and the formal group of a split torus
over $\Z_p$ is torsion-free for $p\ge3$ (\texttt{paper/component-groups.tex}, Lemma~3.1 and Proposition~3.2).
Consequently a torsion CM divisor whose class in $\Phi_p$ vanishes has trivial $\ell$-part for those $\ell$,
which lowers the bound $t$ of Theorem~\ref{thm:counting} for that divisor (Corollary~3.3 there). For an
elliptic quotient $E$ the statement is simpler: a rational torsion point reducing to the singular point lies
outside $E^0$, so when every nonzero torsion point does, the torsion injects into $\Phi_p$
(\texttt{paper/hyperelliptic-15-38.tex}, Lemma~3.1), and the pairing of Proposition~\ref{prop:bielliptic} is
exact as soon as it is torsion and its class in $\Phi_p$ is zero.

Two further points complete the argument in the first shape. The cover $Y\to E$ must be identified among the
double covers branched at the same four points, which are the twists of the hyperelliptic one by $E(\Q)[2]$:
either the dual graph of $Y$ itself decides (the twist by a $2$-torsion class changes the parity of the slope of
the function along the chains of the regular model, and the twisted cover is not semistable where that slope is
odd along a double point of odd thickness; \texttt{paper/hyperelliptic-15-38.tex}, Section~4), or the Prym of
$Y\to E$ does, since the twist has a Prym isogenous over $\Q$ to a product of two elliptic curves while the
$L$-polynomial of the Prym of $Y$ at one good prime may be irreducible of degree four
(\texttt{paper/hyperelliptic-6-385.tex}). And when the branch points are conjugate the pairing lives over a
quadratic field $K$, in the minus part of $E(K)$: the torsion statement is then the Gross--Zagier formula of
Yuan--Zhang--Zhang for the relevant genus character, and the class in $\Phi_p$ can be computed over
$K_{\mathfrak p}$ without matching individual points to components, because a distribution of the points that
would give the conjugate pair a nonzero class is incompatible with the order of the torsion
(\texttt{paper/hyperelliptic-6-323.tex}).

\subsection{Results}\label{sec:candidates-results}

\begin{theorem}\label{thm:thirtythree}
Of the $34$ open quotients $\Curve{D}{N}/W$ on which the differentials satisfy Noether's count for a
hyperelliptic curve, the $33$ listed in Table~\ref{tab:candidates} are hyperelliptic, and
$\Curve{141}{2}/\{1,6,94,141\}$ is not.
\end{theorem}

\todo{At the time of writing, $27$ of the $33$ are proved and $6$ await the dual test of
Proposition~\ref{prop:dual-test} at levels $4620$, $7140$, $6780$, $7788$ and $8260$, where the spaces
$S_{3/2}$ (dimension $561$ to $849$) are being computed; every other step of their proofs is complete and
tested. Update the table and this theorem when they finish.}

\begin{table}[ht]
\centering\small
\renewcommand{\arraystretch}{1.12}
\setlength{\tabcolsep}{4pt}
\begin{tabularx}{\textwidth}{llXl}
\hline
Curve & Shape & How the relation is proved & Note \\
\hline
\multicolumn{4}{l}{\emph{Genus three over genus one (Proposition~\ref{prop:bielliptic})}}\\
$1326_1$, $210_{13}$, $39_{11}$, $55_7$, $6_{301}$, $6_{455}$, $6_{247}$, $6_{253}$, $6_{185}$, $6_{187}$, $22_{51}$, $141_2$, $21_{34}$, $57_{10}$, $15_{46}$, $690_7$
 & 3/1 & exact model: the branch points from their Heegner heights, every cover tested against the trace formula & \texttt{heegner-height-models}\\
$15_{38}$ & 3/1 & dual test; torsion by $\Phi_3,\Phi_5$; twist by the dual graph & \texttt{hyperelliptic-15-38}\\
$6_{385}$ & 3/1 & dual test$^*$; torsion by $\Phi_2$; twist by the Prym & \texttt{hyperelliptic-6-385}\\
$714_5$ & 3/1 & dual test$^*$; torsion by $\Phi_3\times\Phi_7\times\Phi_{17}$; twist by the Prym & \texttt{hyperelliptic-714-5}\\
$6_{323}$ & 3/1 & Gross--Zagier (genus character); torsion by $\Phi_2$ over $\Q(\sqrt{-2})$; twist by the Prym & \texttt{hyperelliptic-6-323}\\
\multicolumn{4}{l}{\emph{Genus four over genus two (Proposition~\ref{prop:normalform})}}\\
$2002_1$ & 4/2 & dual certificate; star torsion-free & \texttt{star-2002-proof}\\
$2030_1$, $2418_1$, $6_{403}$ & 4/2 & three Borcherds forms; star torsion-free & \todo{certificates note}\\
$6_{341}$ & 4/2 & Borcherds forms; the two polarisations in the star's class told apart by the branch divisor & \todo{6\_341 note}\\
$3390_1$ & 4/2 & dual test$^*$; star torsion-free & \texttt{hyperelliptic-3390-1}\\
\multicolumn{4}{l}{\emph{Genus three with the full group (Proposition~\ref{prop:g12})}}\\
$46_{31}$ & 3 & four divisors by Borcherds forms against $t=2$ & \texttt{torsion-counting-hyperelliptic}\\
$6_{379}$ & 3 & dual certificate; torsion-free & \texttt{hyperelliptic-6-379}\\
$21_{58}$ & 3 & dual test at $2\parallel N$; four divisors against $t=2$ & \texttt{hyperelliptic-21-58}\\
$77_{10}$ & 3 & dual test at $2\parallel N$ as a search; five divisors against $t=2$ & \texttt{hyperelliptic-77-10}\\
$39_{34}$ & 3 & dual test at $2\parallel N$, odd fibre via level $N/2$; five divisors against $t=4$ & \texttt{hyperelliptic-39-34}\\
$6_{649}$ & 3 & dual test$^*$; three divisors, $t=6$ lowered to $2$ by $\Phi_3$ & \texttt{hyperelliptic-6-649}\\
$4130_1$ & 3 & dual test$^*$; torsion-free & \texttt{hyperelliptic-4130-1}\\
\hline
\end{tabularx}
\caption{The $33$ hyperelliptic candidates and the proofs. $D_N$ denotes $\Curve{D}{N}/W$ with $W$ as in the
note named; $^*$ marks a dual test whose weight-$\tfrac32$ basis was still being computed at the time of
writing. Equations exist for the first group (Section~\ref{sec:candidates-heights}) and for $2002_1$,
$2030_1$, $2418_1$, $6_{341}$, $6_{403}$ (Proposition~\ref{prop:normalform}); the others are proved
hyperelliptic without one.}
\label{tab:candidates}
\end{table}

\subsection{Equations}\label{sec:candidates-heights}

For a double cover $Y\to E$ of an elliptic curve branched at CM points, the branch points can be located
exactly: the explicit Gross--Zagier formula of Cai, Shu and Tian gives the N\'eron--Tate height of the trace of a
CM point of fundamental discriminant $d$ to $E$, and of its twists by the genus characters, in terms of
$L'(1,E,\chi)$; the points of $E$ of those heights are finitely many, every double cover of $E$ branched along
each candidate divisor is written down, and the trace formula selects the one with the point counts of $Y$
(\texttt{paper/heegner-height-models.tex}, Section~3 and Proposition~2). This produced the twenty equations of
the first group of Table~\ref{tab:candidates}, three of them as double covers of a conic without rational
points, and it is also what shows that $\Curve{141}{2}/\{1,6,94,141\}$ is not hyperelliptic: its unique
compatible cover is a plane quartic. For $2002_1$ the branch fibre of Proposition~\ref{prop:normalform} is
found exactly by a Thue--Mahler equation, since the two branch points stay distinct modulo every prime of good
reduction (\texttt{paper/star-2002-proof.tex}, Proposition~9), and the constant $c$ by the point counts; the
other genus-four models are exact given their stars, whose equations are candidates certified by point counts.
No equation is known for the seven curves of the third group, nor for $15_{38}$, $6_{385}$, $714_5$, $6_{323}$
beyond the models of the heights route, which the proofs above do not use.

\end{document}
